\section{Three boxes}\label{sec:three}

Three boxes require the non-rectangular colour $\yng(2,1)$.
A refined invariant $\rCS_{\yng(2,1)}$ is available by definition for torus knots and iterated torus knots, through the elliptic Hall algebra and the DAHA.
The main results of this section concern these knots.

\subsection{Torus knots and algebraic iterated torus knots}
% [R10, R11]
\begin{compresult}\label{cr:three}
For the $15$ knots of Table~\ref{tab:three}, with the normalisation {\rm(N1)(N2$'$)} of \S\ref{sec:normalization} and at framing zero,
$\hf_\rho$ is positive for all six $\rho$ with $|\rho|\le3$:
with one overall sign for each $\rho$, the spin form $\hbN_{\rho;i,j,k}$ and the genus form $\hbN_{\rho,g,\beta,F}$ consist of non-negative integers.
The list consists of seven torus knots, seven algebraic iterated torus knots, and $\mathbf{4_1}$.
\end{compresult}
\begin{table}[h]
\centering
\caption[Tests up to three boxes at framing zero]{Tests up to three boxes at framing zero. A check mark means that there are no negative terms in the spin form and in the genus form; the number in parentheses is the number of terms of the spin form.
The rows $T_{2,3}$, $T_{2,5}$ and $\mathbf{4_1}$ reproduce tables of KN; for $\mathbf{4_1}$ with $\protect\yng(2,1)$ the KN candidate is used.}\label{tab:three}
\sbox\tabbox{\begin{tabular}{llcccccc}
\hline
Knot & Class & $\yng(1)$ & $\yng(2)$ & $\yng(1,1)$ & $\yng(3)$ & $\yng(2,1)$ & $\yng(1,1,1)$ \\
\hline
$T_{2,3}$ & torus & \checkmark\,(6) & \checkmark\,(26) & \checkmark\,(41) & \checkmark\,(128) & \checkmark\,(236) & \checkmark\,(282) \\ % tells-ok (table entry)
$T_{2,5}$ & torus & \checkmark\,(10) & \checkmark\,(112) & \checkmark\,(144) & \checkmark\,(753) & \checkmark\,(986) & \checkmark\,(1083) \\ % tells-ok (table entry)
$T_{3,4}$ & torus & \checkmark\,(21) & \checkmark\,(429) & \checkmark\,(504) & \checkmark\,(3249) & \checkmark\,(3802) & \checkmark\,(4053) \\ % tells-ok (table entry)
$T_{3,5}$ & torus & \checkmark\,(30) & \checkmark\,(766) & \checkmark\,(867) & \checkmark\,(5831) & \checkmark\,(6559) & \checkmark\,(6890) \\ % tells-ok (table entry)
$T_{2,11}$ & torus & \checkmark\,(22) & \checkmark\,(733) & \checkmark\,(813) & \checkmark\,(5667) & \checkmark\,(6260) & \checkmark\,(6519) \\ % tells-ok (table entry)
$T_{3,7}$ & torus & \checkmark\,(53) & \checkmark\,(1797) & \checkmark\,(1953) & \checkmark\,(13600) & \checkmark\,(14696) & \checkmark\,(15199) \\ % tells-ok (table entry)
$T_{4,5}$ & torus & \checkmark\,(68) & \checkmark\,(2356) & \checkmark\,(2558) & \checkmark\,(18729) & \checkmark\,(20195) & \checkmark\,(20883) \\ % tells-ok (table entry)
$K1$ & algebraic & \checkmark\,(104) & \checkmark\,(4141) & \checkmark\,(4414) & \checkmark\,(32732) & \checkmark\,(34655) & \checkmark\,(35558) \\ % tells-ok (table entry)
$K3$ & algebraic & \checkmark\,(122) & \checkmark\,(5178) & \checkmark\,(5483) & \checkmark\,(41138) & \checkmark\,(43289) & \checkmark\,(44300) \\ % tells-ok (table entry)
$K6a$ & algebraic & \checkmark\,(275) & \checkmark\,(12784) & \checkmark\,(13267) & \checkmark\,(100115) & \checkmark\,(103458) & \checkmark\,(105035) \\ % tells-ok (table entry)
$K9$ & algebraic & \checkmark\,(564) & \checkmark\,(28584) & \checkmark\,(29309) & \checkmark\,(224004) & \checkmark\,(228995) & \checkmark\,(231354) \\ % tells-ok (table entry)
$K8$ & algebraic & \checkmark\,(617) & \checkmark\,(31499) & \checkmark\,(32427) & \checkmark\,(251627) & \checkmark\,(258022) & \checkmark\,(261098) \\ % tells-ok (table entry)
$K7$ & algebraic & \checkmark\,(834) & \checkmark\,(43533) & \checkmark\,(44631) & \checkmark\,(343739) & \checkmark\,(351218) & \checkmark\,(354818) \\ % tells-ok (table entry)
$K10$ & algebraic & \checkmark\,(4119) & \checkmark\,(227727) & \checkmark\,(230642) & \checkmark\,(1805429) & \checkmark\,(1825082) & \checkmark\,(1834640) \\ % tells-ok (table entry)
$\mathbf{4_1}$ & KN & \checkmark\,(10) & \checkmark\,(42) & \checkmark\,(42) & \checkmark\,(223) & \checkmark\,(264) & \checkmark\,(223) \\ % tells-ok (table entry)
\hline
\end{tabular}
}\resizebox{\textwidth}{!}{\usebox\tabbox}
\end{table}
The five torus knots not in the tables of KN ($T_{3,4}$ with three boxes, $T_{3,5}$, $T_{2,11}$, $T_{3,7}$, $T_{4,5}$) are new tables, but since KN conjecture positivity for all torus knots they are further checks of that conjecture.
The three-box results for the seven algebraic iterated torus knots are, to our knowledge, not in the literature.
For the three-strand cable $K7$ and the two-step iteration $K10$ even the refined invariants with $|\lambda|=3$ seem to be new;
Cherednik and Danilenko \cite{Cherednik14} give $\yng(1)$ for $K7$ and the DAHA--Betti polynomial for $K10$.

Since the DAHA and the elliptic Hall algebra determine the values only up to a monomial, the tests for knots not in the KN data depend on the normalisation (N2$'$).
For the cables this rule is Hypothesis~\ref{hyp:N2}, and the dependence is confined to its single point, the colour-dependent power $(q/t)^{c_\lambda}$ not determined at $t=q$.
Shifting the reference framing by an even integer is a shift of $x$ (Proposition~\ref{prop:N2-framing}), and such shifts are covered by the framing scans of \S\ref{sec:three-x}.
Other assignments destroy the Laurent property and the transposition symmetry in the example $K5c$.

% [R12]
The input $\rCS_{\yng(2,1)}$ has negative coefficients, while $\hf_{\yng(2,1)}$ is positive.
For torus knots this was pointed out by KN \cite{Kameyama17} (see also the examples of Cherednik \cite{Cherednik11}); our results show that it persists for algebraic iterated torus knots.
Positivity for three boxes does not follow directly from positivity of the input; it appears only after the extraction of KN (plethystic logarithm and the matrix $M$).

\subsection{Framing scans}\label{sec:three-x}
% [R13]
\begin{compresult}[Direct checks over framings]\label{cr:framing-batch}
In the $1821$ cases (pairs $(x,\rho)$ for a given knot) listed in Table~\ref{tab:framing} there are no negative terms in the spin form or in the genus form; in nine cases, all for the unknot, $\hf=0$.
The largest example is $K10$ with three boxes at $x=2$, with about $1.89$ million terms in the spin form.
\end{compresult}
\begin{table}[h]
\centering
\caption{Framing scans. The columns $|\rho|=1,2,3$ give the range of framings $x$ examined; ``cases'' is the number of pairs $(x,\rho)$ tested (``each'' means per knot;
the case $x=0$ of $K10$ with three boxes belongs to Table~\ref{tab:three} and is not counted here); the last column is the total number of negative terms in the spin form and in the genus form.}\label{tab:framing}
\sbox\tabbox{\begin{tabular}{lccccc}
\hline
Knot & $|\rho|=1$ & $|\rho|=2$ & $|\rho|=3$ & cases & negative (spin/genus) \\
\hline
unknot & $|x|\le6$ & $|x|\le6$ & $|x|\le6$ & 78 & 0/0 \\
$T_{2,3}$ & $|x|\le6$ & $|x|\le6$ & $|x|\le6$ & 78 & 0/0 \\
$T_{2,5}$, $\mathbf{4_1}$ & $|x|\le4$ & $|x|\le4$ & $|x|\le4$ & each 54 & 0/0 \\
$T_{2,7}$, $T_{3,4}$ & $|x|\le3$ & $|x|\le3$ & $|x|\le3$ & each 42 & 0/0 \\
$T_{3,5}$, $T_{2,11}$, $T_{3,7}$, $T_{4,5}$ & $|x|\le10$ & $|x|\le10$ & $|x|\le10$ & each 126 & 0/0 \\
$K1$ & $|x|\le2$ & $|x|\le2$ & $|x|\le2$, $x=\pm6$ & 36 & 0/0 \\
$K3$ & $|x|\le2$ & $|x|\le2$ & $|x|\le2$ & 30 & 0/0 \\
$K6a$, $K7$, $K8$, $K9$ & $|x|\le10$ & $|x|\le10$ & $|x|\le10$ & each 126 & 0/0 \\
$K10$ & $|x|\le10$ & $|x|\le10$ & $|x|\le2$ & 75 & 0/0 \\
$K_p$ ($1\le|p|\le6$) & $|x|\le4$ & $|x|\le4$ & -- & each 27 & 0/0 \\
\hline
\end{tabular}
}\resizebox{\textwidth}{!}{\usebox\tabbox}
\end{table}
In particular, for all knots of Table~\ref{tab:three} the three-box invariants remain positive for $x\neq0$ in the ranges shown.
The torus knot $T_{2,7}$ does not appear in Table~\ref{tab:three}; its test at framing zero for $|\rho|\le3$ is part of the scan of Table~\ref{tab:framing}, which includes $x=0$.
For the cables the scans assume Hypothesis~\ref{hyp:N2}, as does Computational Result~\ref{cr:three}.
The number of terms grows with positive $x$ and decreases with negative $x$ (for $K10$ with $\rho=\yng(3)$, about $1.75$ million terms at $x=-2$ and about $1.86$ million at $x=2$).
For three boxes $\bt^{-3x}\hf^{(x)}_\rho$ is an exponential polynomial with seven bases (Proposition~\ref{prop:expsum}),
but we did not find an ordering of the bases for which the Newton coefficients have a uniform sign, so we have no proof for all $x$.

\subsection{Conjectures}\label{sec:conj}
The evidence of \S\ref{sec:one}--\S\ref{sec:higher} and the picture of \S\ref{sec:physics} lead us to the following conjectures.
% [F19]
\begin{conjecture}[Framed torus knots]\label{conj:A}
For every torus knot $T_{m,n}$, every framing $x\in\Z$ and every partition $\rho$, the framed refined LMOV generating function $\hf^{(x)}_\rho(T_{m,n})$ is positive in the sense of Definition~\ref{def:positivity}:
for one overall sign, the spin form $\hbN^{(x)}_{\rho;i,j,k}$ and the genus form $\hbN^{(x)}_{\rho,g,\beta,F}$ are non-negative integers.
\end{conjecture}
% [R14]
\begin{conjecture}[Algebraic iterated torus knots]\label{conj:B}
For every algebraic iterated torus knot, with the refined invariants defined by the DAHA and normalised by {\rm(N1)(N2$'$)},
$\hf^{(x)}_\rho$ is positive for every framing $x\in\Z$ and every partition $\rho$.
\end{conjecture}
At $x=0$, Conjecture~\ref{conj:A} is part of the conjecture of KN.
We do not formulate a conjecture for general knots with homological input:
for $\rho=\yng(1)$ positivity is a theorem (Corollary~\ref{cor:one}),
for two boxes the knot must satisfy the exponential growth property (Proposition~\ref{prop:two-necessary}), which excludes $\mathbf{9_{42}}$,
and for non-rectangular $\rho$ there is no homological input to speak of.
For the twist and double twist knots, Theorems~\ref{thm:twist}--\ref{thm:twist-x} support positivity of two boxes at all framings, given the conjectural input.

The motivation for Conjecture~\ref{conj:B} is different from that for torus knots.
The physical argument of KN obtains the additional $U(1)_R$ symmetry from the Seifert structure.
The complement of an algebraic iterated torus knot is in general not a Seifert manifold, so the argument does not apply directly.
The motivation comes instead from two conjectures of Cherednik and Danilenko \cite{Cherednik14}.
Their positivity conjecture states that for algebraic knots and rectangular Young diagrams all coefficients of the DAHA superpolynomial are positive.
Their Connection Conjecture (ii) states that for uncoloured knots with this positivity property (called pseudo-algebraic) the DAHA superpolynomial coincides with the tilde-normalised reduced stable Khovanov--Rozansky polynomial.

For $\yng(1)$, positivity follows from the Connection Conjecture (ii) and Proposition~\ref{prop:one}.
Indeed, under this conjecture the DAHA value for $\yng(1)$ is a monomial multiple $m\,P_{\mathrm{KhR}}$ of the Poincar\'e polynomial of Khovanov--Rozansky homology.
By (N1) both give the same reduced HOMFLY-PT polynomial at $\bt=-1$, so $m=\pm\bt^k$ with $\pm(-1)^k=1$; since both have non-negative coefficients, $m=\bt^{2l}$.
As $\bt^{2l}$ is $\sigma$-invariant, Proposition~\ref{prop:one} applies and positivity for $\yng(1)$ follows.
For rectangular colours, Cherednik and Danilenko conjecture positivity of the DAHA values and expect agreement with homology; the rectangular part of Conjecture~\ref{conj:B} transfers this to the level of BPS invariants.
For $\yng(2,1)$, on the other hand, Cherednik and Danilenko make no positivity claim, so that positivity at the level of BPS invariants lies outside their conjectures.
As for the restriction to algebraic knots, Cherednik and Danilenko note that there are non-algebraic iterated torus knots without positivity,
and for most of the non-algebraic cables that we computed, negative coefficients appear already in the DAHA value for $\yng(1)$ (Appendix~\ref{app:nonalg}).

\subsection{Non-rectangular colours for hyperbolic knots}\label{sec:21}
For hyperbolic knots there is no definition-based refined invariant for $\yng(2,1)$, either in the literature or in software.
We examined whether structural conditions can determine it.
% [R15, R16, R17]

Working knot by knot, we imposed the coloured HOMFLY-PT polynomial for $\yng(2,1)$ at $t=q$ (obtained from three-strand cables), the mirror symmetry, exponential growth, the reduction to $U(2)$,
the Laurent property of the three-box $\hf$, and the closure of the genus expansion.
These conditions do not determine $\rCS_{\yng(2,1)}$ uniquely: the difference of two solutions has the form of a knot-independent factor times a free polynomial.
The existence of an ambiguity was already noted by Anokhina et al.\ \cite{Anokhina12}.
One can still ask whether the conditions admit a solution compatible with positivity.
An integral $\rCS_{\yng(2,1)}$ satisfying positivity exists for $\mathbf{4_1}$ (the KN candidate) and for $\mathbf{6_1}$ (a solution with $813$ terms).
For $\mathbf{5_2}$ we found only rational solutions; the existence of an integral solution is undecided.
This is not a test of the conjecture, only a check that it is not contradicted.

We also considered families of knots.
At $t=q$, the formula of Morozov \cite{Morozov16,Morozov19} for double twist knots coloured by $\yng(2,1)$ agrees with the values from cables; this is a check of a known formula.
For refined rectangular colours we implemented the KNTZ formulae and compared them with the GNSSS formulae and with the elliptic Hall algebra values for the trefoil up to $\yng(2,2)$.
The refined $\yng(2,1)$, however, is not determined uniquely even under assumptions on the structure of the family.
Three boxes for hyperbolic knots therefore require a definition of the refined invariant for $\yng(2,1)$; candidates are discussed in \S\ref{sec:discussion}.
